\chapter{结论与展望}

% TODO(P-002): This chapter mirrors Section 1.4. Do not introduce claims unsupported by preceding chapters.

\section{全文工作总结}

三维医学图像分割需要在统一体数据空间中组织多种层次的信息。一方面，器官和组织的边缘、纹理及邻域结构依赖局部空间表征，跨切片、跨区域的解剖联系又要求模型能够建立较大范围的上下文关系；另一方面，U形网络在逐级编码与解码过程中持续改变特征的空间分辨率、信息汇聚范围和语义状态，使同一结构在不同尺度层级中的表达条件并不相同。围绕这一总体问题，本文没有将长距离建模与多尺度建模作为彼此孤立的技术目标，而是形成由空间关系到尺度层级关系的递进研究路线：先处理局部结构信息与长距离空间上下文之间的协同，再在已经形成的局部--全局表征框架上研究不同尺度结构信息如何保持层级对应并参与解码重建，由此将局部结构、长距离上下文和尺度层级结构信息纳入同一三维分割信息组织主线。

针对空间维度的信息协同问题，本文构建了基于Vision-xLSTM局部--全局协同建模的三维医学图像分割方法ViL-UNet。该方法以3D CNN作为局部与层次特征编码基础，使体数据中的邻域模式和多层空间表示首先在规则三维网格中形成；随后将深层三维特征组织为视觉序列，利用Vision-xLSTM建立较长距离的上下文交互，并通过多方向序列组织扩展状态更新对不同空间方向的访问关系。解码路径从深层上下文表示逐级恢复空间分辨率，同时结合跨层特征重新组织局部细节与高层语义，使卷积形成的局部结构信息和序列建模获得的长距离上下文共同进入体素级预测。由此，第一项研究并非以序列模型替换卷积编码，而是使两类建模方式分别承担局部空间表征和长距离关系组织功能，在统一U形框架中形成连续的局部--全局协同信息流。

围绕上述方法，本文在Synapse腹部CT和ACDC心脏MRI数据集上完成了总体性能、器官或结构级表现、ViL深度配置、模型复杂度和可视化等多层次分析。在Synapse数据集上，ViL-UNet的平均DSC为85.12\%，HD95为12.49~mm，在所列比较方法中分别对应最高平均DSC和最低HD95；在ACDC数据集上，其平均DSC达到92.79\%，同样位于表内最高位置。两套数据覆盖不同成像模态和解剖区域，相关结果从区域重叠和边界距离等角度支持了局部空间表征与长距离上下文协同的整体设计。消融与复杂度分析还表明，ViL模块数量增加与性能变化并非简单单调关系，最终采用的配置与所测试深度范围内的结果相一致；同时，模型在保持较高分割结果的情况下具有较为紧凑的参数规模。第一项研究由此建立了后续尺度结构研究所依赖的局部--全局表征基础。

在空间关系得到组织之后，本文进一步将问题推进到U形网络内部的尺度维度。不同编码层级具有不同空间分辨率和语义状态，扩大空间交互范围并不会自动解决各尺度结构信息应如何形成、如何与当前解码状态建立联系以及如何参与空间恢复的问题。第四章的前置对比显示，持续累积历史尺度信息后，不同评价维度没有形成方向一致的有利变化，这一现象将设计重点由“增加跨尺度传播”转向“保持尺度对应并按当前层级状态利用结构信息”。基于此，本文构建基于尺度特异结构证据调制的三维医学图像分割方法：各尺度分别依据自身层次特征形成对应的结构证据，使结构信息来源保持尺度独立；当前解码状态与对应结构证据共同确定通道级调制关系，使结构信息的利用程度与当前重建状态建立联系；经联合映射形成的结构增强分量再以残差形式接入解码表示，在保留基础解码直接路径的同时，使尺度结构信息参与相应层级的空间恢复。整个过程形成“当前尺度特征—尺度特异结构证据—当前状态条件下的自适应调制—残差增强—对应尺度重建”的连续信息链。

第四章围绕这一尺度信息组织链，从总体性能、器官级差异、尺度传播策略、不同结构信息组织配置、模型复杂度与运行开销以及困难样本等方面展开分析。总体比较联合考察区域重叠与边界距离，器官级结果呈现出不同解剖目标之间并不完全一致的变化模式；尺度传播策略比较将历史累积问题与尺度相对独立的信息组织方式联系起来。不同配置的对比还呈现出非单调的中间结果：建立独立尺度结构证据及其解码接入路径后，仅把自适应调制限制在单一解码层级并不会自动形成更优结果，完整方法因而更强调多个对应解码阶段之间协调的结构信息利用。与此同时，复杂度分析将性能变化与新增计算、显存和推理开销置于同一评价框架，困难样本讨论则关注局部连续性、小结构和复杂边界等空间现象及仍可能存在的误差。由此，本文两层研究在同一三维分割框架中形成明确递进：空间维度解决局部结构与长距离上下文的协同，尺度维度进一步组织不同层级的结构信息及其与当前解码状态的关系，最终形成从局部结构建模、长距离空间上下文组织到尺度层级结构信息利用的统一技术主线。

\section{主要研究结论}

围绕空间维度的信息建模，可以得到的第一点认识是，局部结构表征与长距离上下文之间具有互补关系。三维医学影像中的边缘、纹理、邻域连续性等信息依赖对局部空间模式的稳定刻画，而器官内部及不同区域之间的空间联系又需要更大范围的上下文参与判断。第一项研究在两个公开基准数据集上的实验结果支持了将两类信息按照不同功能组织在同一分割框架中的合理性：局部建模为细粒度空间表征提供基础，长距离关系建模则补充跨区域、跨切片的信息交互。由此，三维分割中的空间上下文建模不宜被理解为对局部表征的替代，而应关注不同关系跨度的信息如何协同服务于同一体素级预测过程。局部信息决定模型能否保留细粒度的空间差异，较大范围上下文则为这些局部响应提供更完整的位置与结构关系参照；只有明确两类信息的功能分工，扩大建模范围才不会以削弱局部表达为代价。

围绕尺度维度的信息组织，可以进一步得到第二点认识：跨尺度信息利用的关键并不在于持续增加历史传播量，而在于结构信息能否保持与当前层级的明确对应，并根据当前解码状态得到适当利用。前置对比显示，当历史尺度状态持续累积时，不同评价维度并未形成方向一致的有利变化，这说明更多的跨尺度历史信息并不等同于更合理的尺度结构组织。对于具有层级编码与解码过程的三维分割网络，更值得关注的是结构信息从何种尺度表征中形成、不同尺度之间是否保留相对独立的组织关系，以及对应结构信息如何与当前层级已经形成的重建状态建立联系。尺度建模因而应从单纯扩大跨层传播转向保持来源对应、层级对应与使用条件之间的一致关系。

综合上述两层研究，可以得到全文层面的第三点认识：三维医学图像分割中的信息组织需要区分问题所发生的维度以及不同信息在该维度中承担的功能。空间维度关注不同位置之间的信息交互范围，核心在于局部结构与长距离上下文的协调；尺度维度关注不同分辨率和语义层级中的结构信息如何形成并参与重建，核心在于尺度特异结构信息与当前解码状态之间的对应关系。因此，三维分割中的信息建模不应简单归结为扩大上下文范围或增加多尺度融合来源，而应根据空间关系和层级关系分别确定信息来源、交互对象与作用方式，并在统一分割框架中建立相互衔接的组织关系。这一区分也意味着，设计新的信息交互机制时需要先明确其要处理的是空间跨度问题还是层级尺度问题，再判断新增信息是否与当前表征状态具有合理对应，而不能仅以信息数量或连接数量衡量建模充分性。相应地，对方法的评价也需要同时关注区域重叠、边界质量以及不同结构上的响应差异，而不能由单一总体指标替代对信息组织效果的完整判断。

\section{局限性与展望}

% TODO: Derive future work from verified limitations, difficult cases, or failure evidence.
